\chapter{Conclusion and open problems}\label{ch-conclusion}\label{sec-further}

\section{Summary of what is proved}
The thesis establishes four results, listed here in the order in which they
are used rather than in order of importance.

\begin{enumerate}[leftmargin=2.2em,itemsep=4pt]
\item \textbf{Attainment and two-sided bracketing}
(Lemma~\ref{attainment}). The variational quantity $\mathcal M=\sup_W
\Sigma_3(W)$ is attained, every maximiser has three positive leading
eigenvalues, and $\cstar\le\mathcal M\le(1+\sqrt3)/2$ with
$\cstar=(9+\sqrt5+2\sqrt{35})/18$ produced by an explicit six-part step
graphon modelled on $K_1\vee C_5$.

\item \textbf{Finite canonical support} (Theorem~\ref{structure-main}). The
canonical law of \emph{every} maximiser is a convex combination of at most
$24$ points of the ellipsoid $\{v^TLv=\mathcal M\}$; the maximiser is almost
everywhere the binary threshold kernel determined by the signs of the pairwise
dot products; the support points satisfy the exact algebraic system
$\sum_ip_iv_iv_i^T=I_3$, $v_i^TLv_i=\mathcal M$, $\sum_jp_jB_{ij}v_j=Lv_i$;
and every positive clique has at most six linearly independent types.

\item \textbf{Joint moment certificate} (Theorem~\ref{main}). Every Borel
probability law on $\R^3$ with $\E XX^T=I_3$ satisfies
$\E|X\cdot X'|+|\E X|^2\le2c$ with
$c=1.284673255723918505\ldots$. The proof is an exact identity in which the
mean deficit, the radial variance, the regression error, the two harmonic
moments, the quartic slack and the atomic self-pair term all appear with
explicit nonnegative coefficients.

\item \textbf{Finite bound and its scope} (Theorem~\ref{main},
Propositions~\ref{weighted-form} and~\ref{no-gain}). Every graph of order
$n\ge3$ satisfies $\Lambda_3(G)\le cn-3$; the same inequality holds for every
symmetric matrix with zero diagonal and off-diagonal entries in $[0,1]$; and
passing from graphs to such matrices does not change the asymptotic extremal
constant.
\end{enumerate}

The improvement over~\cite[Proposition 6.1]{STW} is a coefficient decrease
of $0.02434374\ldots$ with the constant term $-3$ unchanged; the improvement
over~\cite[Theorem 3.4.1]{Ebrahimi} is $3/\sqrt5-c=0.05696753\ldots$. Both
statements are proved, not numerical, and both constants are printed to sixty
digits in Section~\ref{sec-protocol}.

\section{The residual gap}
% Ported verbatim from paper/top3_joint_moment_bound.tex lines 1386-1414.
% Only cross-reference wording was adapted; no mathematics was altered.
The explicit construction in Lemma~\ref{attainment} leaves the interval
\[
 \frac{9+\sqrt5+2\sqrt{35}}{18}\le\mathcal M
       \le1.284673255723918505\ldots.
\]
The lower endpoint comes from a weighted wheel, but its global optimality
is not established here. Its provenance is the subject of
Remark~\ref{rem-newton}: the same constant and the same optimising weight
were announced earlier, conditionally, in~\cite{Newton1,Newton2}, whose
certificate argument would close the interval at the lower endpoint were it
to survive the independent audit that its author requests. Independently of
that outcome, the unconditional question left open here is the gap
\[
 c-\cstar=0.003105058851750628\ldots,
\]
and removing it is the natural next problem. The support bound reduces the
structure of every
maximizer to finitely many positive-mass types without classifying the
possible templates. It does not imply that the binary matrix $B$ is
nonsingular, that every maximizer has six types, or that individual masses
are uniformly bounded away from zero. Determining the optimum remains
a separate problem.

The exact identity~\eqref{certificate-final} also distinguishes a
quantitative certificate from a qualitative equality obstruction. Any
further decrease obtained from this certificate must bound a remaining term from below
under the same constraints, without counting the same quartic remainder
twice. The present argument does not use an unquantified stability gap.


\section{Open problems}\label{sec-open}
The following are stated in decreasing order of how much of the present
machinery they would reuse. None is proved here.

\begin{enumerate}[leftmargin=2.2em,itemsep=5pt]
\item \textbf{Close the gap.} Prove or disprove $\mathcal M=\cstar$, i.e.\
remove the gap $c-\cstar=0.00310506\ldots$. The obstruction is identified
in Appendix~\ref{sec-appendix}: the certificate loses through the non-sharp
constant $3/2$ of Lemma~\ref{contact-square} (the true supremum is
$1.200224\ldots$ on $[0,1]$), through the divisor $125$ in~\eqref{Cstar8},
and through the endpoint-ratio hypotheses $3$ and $\sqrt8$ that are stronger
than needed. Replacing $3/2$ by a sharp contact constant is the cheapest
identified improvement and is a finite exact-arithmetic computation, but it
was not carried out.

\item \textbf{Determine the joint feasible region.} Let
$\Gamma=\overline{\{(\E|X\cdot X'|,|\E X|^2)\}}$ over isotropic laws on
$\R^3$. The two coordinate bounds are known and sharp,
$\E|X\cdot X'|\le(1+\sqrt5)/2$ and $|\E X|^2\le1$, and Theorem~\ref{main}
bounds their sum.
The region $\Gamma$ itself is not known; it is the natural two-dimensional
formulation of the mixed projection constant raised in~\cite[Section 6.2]{STW}.
This is the problem for which the present method gives a supporting line
rather than the whole boundary.

\item \textbf{Classify the extremal templates.} The support theorem bounds the
number of types but does not classify them. In particular it is not known
whether $B$ is nonsingular, whether every maximiser has exactly six types, or
whether the masses are bounded away from zero uniformly over maximisers. The
numerical evidence of Chapter~\ref{ch-numerical} points to the family
$K_m\vee C_5[\cdot]$ but proves nothing.

\item \textbf{Exact finite-order extremal graphs.} Determine $M_3(n)$ and the
graphs attaining it for every $n$. At present this is known only for $n\le7$
by exhaustive computation (Chapter~\ref{ch-numerical}); the extremal graphs
there are $K_1\vee C_5$ for $n=6$ and $K_2\vee C_5$ for $n=7$, which suggests
a blow-up family with a congruence-dependent description as in~\cite{HW}, but
no conjecture strong enough to state is supported by the data.

\item \textbf{Resolve the status of the announced certificate.} The
deposits~\cite{Newton1,Newton2} assert $\Lambda_3(G)\le\cstar n-3$
conditionally. Their author requests an independent audit of the admissibility
of the archived dual payments, the global sector partition and the exact
verifier. Completing that audit would either collapse the interval of
Section~\ref{sec-further} to a point or produce a counterexample to a
published-style claim.

\item \textbf{Ring length.} The construction uses $C_5$. Whether
$K_1\vee C_r$ for $r\ne5$, or more general joins, produce larger constants has
not been investigated here and is recorded as an open computation rather than
as a conjecture.
\end{enumerate}

\section{Adjacent problems}\label{sec-adjacent}
Three neighbouring spectral-sum problems frame the present one and are worth
stating precisely, because the methods differ.

\emph{The two-term sums.} $\Lambda_2=\lambda_1+\lambda_2$ is settled:
$\Lambda_2(G)\le8n/7-2$ with equality exactly when $7\mid n$, and the unique
extremal graph is the three-part join described in~\cite{HW}, building
on~\cite{KLMPT}. For $\lambda_2+\lambda_3$ the bound $2n/3-2$ follows from the
explicit description of the rescaled pair region in~\cite{Paulin}. A separate
line of the present author's work treated $\lambda_2+\lambda_3$ with an
equality characterisation and a finite-order stability refinement; it was
completed independently of and concurrently with~\cite{Paulin} and is not
reproduced here, since its subject is a different functional.

\emph{Weighted two-term sums.} For $\lambda_1+k\lambda_2$ the asymptotic
coefficient is already determined for every fixed $k\ge0$ by the support
function of the region in~\cite{Paulin}:
\[
 h(k)=\frac{k+3+2\sqrt{2k^2-2k+4}}7\ \ (0\le k\le3),
 \qquad h(k)=\frac{k+1}2\ \ (k\ge3),
\]
so that $\lambda_1+k\lambda_2\le h(k)n-(k+1)$ and the coefficient is
asymptotically sharp. For $k=2$ this is
$h(2)=(5+4\sqrt2)/7=1.5224077499\ldots$. What remains open for
$\lambda_1+2\lambda_2$ is exactly what remains open for $\Lambda_3$: the exact
finite-order maximum, the full list of extremal graphs, and a stability
theorem. For trees and large order, the extremal tree is known~\cite{KMMZ}.

\emph{Individual eigenvalues.} $\lambda_3(G)\le n/3-1$ is sharp
by~\cite{STW}, and $c_3$ is therefore determined; the interesting individual
problem is now $c_k$ for $k\notin\{2,3,4,8,24\}$, where the upper bound
$\alpha_k$ of~\cite{STW} and the lower bounds of~\cite{Linz} do not meet.

\section{What would count as a complete solution}
For the reader who wants a single criterion: the problem posed in
Section~\ref{sec-problem} is solved when one exhibits an explicit constant
$c^\dagger$ and an explicit infinite family of graphs with
$\Lambda_3(G)=c^\dagger|V(G)|-3+o(1)$, together with a proof that
$\Lambda_3(G)\le c^\dagger|V(G)|-3$ for all $G$. At present $c$ is proved and
$\cstar$ is attained asymptotically by the weighted blow-ups of $K_1\vee C_5$,
but $\cstar$ is not proved and $c$ is not known to be attained. Closing either
half would be a substantial result; closing both at the same value would
finish the problem.
